\section{Traducción automática no supervisada: la intuición geométrica solo es el punto de arranque}

La pregunta de investigación temprana de Lample era muy exigente: si solo se cuenta con textos monolingües de cada uno de dos idiomas, sin un corpus paralelo a nivel de oración, ¿es posible aprender a traducir? En la exposición oral de la clase se describe el hallazgo central de forma intuitiva: los espacios de vectores de palabras de distintos idiomas pueden tener estructuras geométricas similares y, tras una transformación, pueden alinearse. Esta intuición es importante, pero si solo se escribe como «los dos espacios difieren en una rotation», se omite el ciclo de entrenamiento que realmente hace funcionar al sistema.

\subsection{Del espacio de representación al diccionario bilingüe}

Sea $X\in\mathbb{R}^{d\times n}$ la matriz de vectores de palabras del idioma de origen y $Y\in\mathbb{R}^{d\times n}$ la matriz de puntos de anclaje del idioma de destino. Un paso de alineación clásico puede escribirse como el problema ortogonal de Procrustes:
\[
W^{\star}=\arg\min_{W}\|WX-Y\|_{F}^{2},
\qquad \text{s.t.}\quad W^{\top}W=I.
\]
Aquí $W$ es una transformación ortogonal que conserva longitudes y ángulos, y $\|\cdot\|_{F}$ es la Frobenius norm. La intuición detrás de la restricción $W^{\top}W=I$ es la siguiente: si la geometría semántica local de los dos idiomas es similar, la alineación debe lograrse en lo posible mediante rotaciones o reflexiones, y no estirando el espacio de forma arbitraria. Un sistema real aún debe encontrar primero puntos de anclaje aproximados, tratar la frecuencia de las palabras y la polisemia, y enfrentar el hecho de que las estructuras de los dos idiomas no son completamente isomorfas.

La siguiente figura separa el «arranque geométrico» del «algoritmo de aprendizaje completo». Primero se infiere un diccionario aproximado a partir de representaciones monolingües y después se mejoran repetidamente el codificador y el decodificador mediante denoising y back-translation.

\lecturefigure{02-cross-lingual-alignment.png}{La traducción no supervisada arranca con la alineación de representaciones y luego mejora de forma gradual mediante entrenamiento cíclico.}{Subtítulos locales 03:00--05:00; Lample et al., \emph{Unsupervised Machine Translation Using Monolingual Corpora Only}, 2018.}

\begin{knowledgebox}{Lectura de la figura: por qué la rotation no puede resolver la traducción por sí sola}
Los vectores de palabras alineados, como mucho, aportan candidatos a nivel de palabra; no resuelven por sí mismos el orden de las palabras, los cambios morfológicos, las dependencias de largo alcance ni la ambigüedad contextual. El sistema del artículo también necesita denoising auto-encoding para que el modelo aprenda la estructura interna de cada idioma, y back-translation para traducir de nuevo la traducción actual al idioma original y formar así una señal autosupervisada. La similitud geométrica responde «por dónde empezar»; el entrenamiento cíclico responde «cómo seguir mejorando».
\end{knowledgebox}

\begin{warningbox}{No confundir la similitud de representaciones con un isomorfismo completo entre idiomas}
Que los espacios de vectores de palabras sean aproximadamente alineables depende del corpus, la frecuencia de las palabras, el método de entrenamiento y la distancia entre los idiomas. En idiomas con pocos recursos, idiomas de morfología compleja o corpus con grandes diferencias de dominio, el diccionario aproximado puede presentar sesgos sistemáticos. Lo que este método demuestra es que «es posible construir una señal de arranque no supervisada», no que «cualquier par de idiomas difiere solo en una matriz».
\end{warningbox}

\teachervoice{Aquí el docente conserva un hábito de investigación muy valioso: buscar primero la estructura de baja dimensión en un problema de alta dimensión. Los vectores de palabras parecen difíciles de observar directamente, pero si se supone que los dos espacios comparten relaciones geométricas, un problema sin etiquetas puede convertirse en un problema de alineación optimizable. Después, el ciclo de entrenamiento corrige esta suposición demasiado fuerte.}

\subsection{Resumen de la sección}

La traducción automática no supervisada muestra que un buen sistema de investigación suele estar compuesto por una «suposición de arranque interpretable» y un «ciclo de retroalimentación capaz de corregirse a sí mismo». Quedarse solo con la rotation sobrestima la suposición; ignorar la back-translation subestima la ingeniería del sistema.
